
\documentclass{article}
\usepackage{colortbl}
\usepackage{makecell}
\usepackage{multirow}
\usepackage{supertabular}

\begin{document}

\newcounter{utterance}

\centering \large Interaction Transcript for game `cladder', experiment `full\_v1.5\_default', episode 8201 with selfrefine{-}Qwen3.5{-}9B.
\vspace{24pt}

{ \footnotesize  \setcounter{utterance}{1}
\setlength{\tabcolsep}{0pt}
\begin{supertabular}{c@{$\;$}|p{.15\linewidth}@{}p{.15\linewidth}p{.15\linewidth}p{.15\linewidth}p{.15\linewidth}p{.15\linewidth}}
    \# & \multicolumn{2}{c}{Player} && \multicolumn{2}{c}{Game Master} \\
    \hline

    \theutterance \stepcounter{utterance}  
    & & & \multicolumn{4}{p{0.6\linewidth}}{
        \cellcolor[rgb]{0.9,0.9,0.9}{
            \makecell[{{p{\linewidth}}}]{
                \texttt{\tiny{[P1$\langle$GM]}}
                \texttt{You are an expert in causal inference. The following question is not a typical commonsense query, but rather a meticulously designed question created by a professor specializing in causal inference, intended to assess the students' mastery of the course content.} \\
\texttt{Imagine a self{-}contained, hypothetical world with only the following conditions, and without any unmentioned factors or causal relationships: Tijv has a direct effect on tijw and xevo. Xevo has a direct effect on gyzp. Tijw has a direct effect on gyzp. For those who are not xevo and are not tijw, the probability of gyzp is 76\%. For those who are not xevo and are tijw, the probability of gyzp is 76\%. For those who are xevo and are not tijw, the probability of gyzp is 12\%. For those who are xevo and are tijw, the probability of gyzp is 14\%. For those who are not xevo, the probability of tijw is 28\%. For those who are xevo, the probability of tijw is 14\%. For those who are xevo, would it be more likely to see gyzp if the individual was not xevo?} \\
\texttt{Start your answer with 'Yes' or 'No', followed by additional reasoning or evidence to support your explanation.} \\
            }
        }
    }
     \\ \\

    \theutterance \stepcounter{utterance}  
    & \multicolumn{4}{p{0.6\linewidth}}{
        \cellcolor[rgb]{0.561,0.737,0.561}{
            \makecell[{{p{\linewidth}}}]{
                \texttt{\tiny{[P1$\rangle$GM]}}
                \texttt{problem} \\
\texttt{The next planned game move requires refinement to adhere to the game's specific formatting and strategic constraints.} \\
\\ 
\texttt{**Format Violation:** The move begins with "No" followed by a detailed derivation. While the prompt asks to start with 'Yes' or 'No', the subsequent text in the "Reasoning" section of the *feedback* explicitly states that the draft "fails the format requirement by not starting with 'Yes' or 'No'". However, looking closely at the draft provided in the "next planned game move" section, it *does* start with "No". The feedback seems to be misinterpreting the draft or referring to a different version. Regardless, the feedback also claims the move "violates the game's constraint to not produce a new answer by providing a full derivation instead of evaluating the draft". This is a contradiction: the user is asking me to *refine* the move, which inherently involves producing a new version (a new answer). The feedback's instruction to "not produce a new answer" is likely a constraint on the *AI's role* in a different context (e.g., a critique{-}only mode), but here the user explicitly asks to "refine this move".} \\
\\ 
\texttt{However, the most critical issue identified in the feedback is the strategic one: "incorrectly judges the draft's strategy as problematic when the draft's reasoning is actually sound despite the formatting error."} \\
\\ 
\texttt{Let's re{-}evaluate the draft's reasoning:} \\
\texttt{1.  **Question:** For those who are \$xevo\$, would it be more likely to see \$gyzp\$ if the individual was not \$xevo\$?} \\
\texttt{    *   This asks to compare \$P(gyzp \textbackslash{}mid xevo=1)\$ vs \$P(gyzp \textbackslash{}mid do(xevo=0))\$.} \\
\texttt{2.  **Draft Calculation for \$P(gyzp \textbackslash{}mid xevo=1)\$:**} \\
\texttt{    *   \$P(tijw=1 \textbackslash{}mid xevo=1) = 0.14\$.} \\
\texttt{    *   \$P(gyzp \textbackslash{}mid xevo=1, tijw=1) = 0.14\$.} \\
\texttt{    *   \$P(gyzp \textbackslash{}mid xevo=1, tijw=0) = 0.12\$.} \\
\texttt{    *   Total = \$0.14 \textbackslash{}times 0.14 + 0.86 \textbackslash{}times 0.12 = 0.0196 + 0.1032 = 0.1228\$.} \\
\texttt{    *   This calculation is correct based on the provided text.} \\
\texttt{3.  **Draft Calculation for \$P(gyzp \textbackslash{}mid do(xevo=0))\$:**} \\
\texttt{    *   The draft cuts off here: "If we set \$xevo = 0\$, the direct effect of \$xevo\$ on \$gyzp\$ is removed. However} \\
            }
        }
    }
    & & \\ \\

    \theutterance \stepcounter{utterance}  
    & & & \multicolumn{2}{p{0.3\linewidth}}{
        \cellcolor[rgb]{0.9,0.9,0.9}{
            \makecell[{{p{\linewidth}}}]{
                \texttt{\tiny{[GM$|$GM]}}
                \texttt{ParseError: The response doesn't start with one of the required word {[}'yes', 'no'{]}, but problem} \\
            }
        }
    }
    & & \\ \\

    \theutterance \stepcounter{utterance}  
    & & & \multicolumn{2}{p{0.3\linewidth}}{
        \cellcolor[rgb]{0.9,0.9,0.9}{
            \makecell[{{p{\linewidth}}}]{
                \texttt{\tiny{[GM$|$GM]}}
                \texttt{game\_result = ABORT} \\
            }
        }
    }
    & & \\ \\

\end{supertabular}
}

\end{document}
